\aufzaehlungfuenf{Die Abbildung
\maabbeledisp {\pi_{\text{vert} }} {TE= TX \times \R \times \R } { p^*E  = X \times \R \times \R
} {(P,u,v,w)} { (P,u,  \omega(P,v) + w)
} {}
ist bei fixiertem $(P,u)$ linear in
\mathkor {} {v} {und} {w} {,}
es liegt also ein Homomorphismus von Vektorbündeln vor. Die Abhängigkeit von der Basis $E$ ist stetig differenzierbar, da $\omega$ stetig differenzierbar ist. Es liegt eine Spaltung zur natürlichen Inklusion
\maabb {} {p^*E} { TE
} {}
vor, da
\mathl{(P,u,w)}{} auf
\mathl{(P,u,0,w)}{} und dieses auf

\mavergleichskette
{\vergleichskette
{ (P,u, \omega(P,0) + w)
}
{ = }{ (P,u,  w)
}
{  }{
}
{    }{
}
{   }{
}
}
{}{}{}
abgebildet wird.
}{Zu
\maabb {f} { U } {\R
} {}
müssen wir die Hintereinanderschaltung

\mathdisp {TX \stackrel{T(f)}{  \longrightarrow } T E = TX \times \R \times \R \stackrel{ \pi_{\text{vert} } }{ \longrightarrow } X \times \R \longrightarrow \R} {  }

betrachten, und diese ergibt

\mathdisp {(P,v) \longmapsto  (P,f(P), v, (df)_P(v) ) \longmapsto (P, f(P), \omega(P,v) + (df)_P(v) ) \longmapsto  \omega(P,v) + (df)_P(v)} { . }

}{Ein horizontaler Schnitt liegt genau dann vor, wenn die vertikale Ableitung verschwindet, und dies bedeutet hier

\mavergleichskettedisp
{\vergleichskette
{ \omega(P,v) + (df)_P(v)
}
{ =} { 0
}
{ } {
}
{ } {
}
{ } {
}
}
{}{}{}
für alle

\mavergleichskette
{\vergleichskette
{ (P,v)
}
{ \in }{ TX
}
{  }{
}
{    }{
}
{   }{
}
}
{}{}{.}
}{Dies beruht darauf, dass eine Form genau dann exakt ist, wenn sie lokal eine Stammform besitzt.
}{Dies folgt  aus dem bisher Bewiesenen.
}

 [[Kategorie:Latexseite]]
